\chapter{General Conclusion}
\label{ch:conclusion}

% ============================================================
\section{Summary of the Dissertation}
\label{sec:conclusion_summary}
% ============================================================

This dissertation examined machine learning for two distinct \gls{ALS} prognosis outcomes using the \gls{PRO-ACT} database. Study~I addressed short-term functional progression, whereas Study~II considered survival within 24 months of symptom onset. The targets and experimental designs were different, but both studies tested how outcome construction, class imbalance, validation, threshold choice, and post-hoc explanation affect the interpretation of performance.

Study~I developed the dissertation's original modelling pipeline. Thirty-five baseline features were used to classify rapid and slow \gls{ALSFRS-R} progression at three- and six-month horizons. Participants were separated before model development, target cutoffs were estimated from training data, and seven classifier families were compared using \gls{PR-AUC}. The selected XGBoost model was then evaluated on the reserved six-month test partition and examined using \gls{SHAP}, \gls{LIME}, Logistic Regression coefficients, subgroup analyses, and a prototype patient-level report.

Study~II moved from original model development to partial replication. It reconstructed the fixed-horizon survival task and 23-feature representation reported by Papaiz et al.~\cite{Papaiz2024}, re-examined the effect of BalancedBagging, and extended the comparison to LightGBM and ten imbalance-handling strategies. \gls{ROC-AUC} replaced the reference study's balanced accuracy as the optimisation criterion. Held-out sensitivity, specificity, balanced accuracy, uncertainty intervals, and alternative thresholds were then needed to interpret the selected models.

The shared contribution is therefore methodological rather than a single prognostic model: both studies link explicit outcome construction and separated evaluation with uncertainty estimates and cautious interpretation of post-hoc explanations.

% ============================================================
\section{Main Findings from Study~I}
\label{sec:conclusion_study1}
% ============================================================

The primary six-month cohort contained 1,392 participants. Among 13 tuned model--mode configurations, unbalanced XGBoost ranked first on development \gls{PR-AUC} ($0.434 \pm 0.049$), narrowly ahead of unbalanced Random Forest ($0.431 \pm 0.042$). The fold distributions overlapped and the Friedman test did not reject equal model ranks ($p=0.053$), so XGBoost was selected by the prespecified ranking rule rather than by evidence of statistical superiority. The three-month sensitivity analysis produced a similar range of cross-validated \gls{PR-AUC} values, although the eligible cohort, slope estimates, and exact ordering differed.

Repeated stratified validation of the saved configurations changes that ordering: balanced Random Forest leads at six months ($0.422\pm0.034$), while unbalanced XGBoost obtains $0.414\pm0.033$. No new tuning or test-informed replacement followed. The result makes the original narrow lead less decisive and shows why the selected model should not be described as intrinsically superior.

On the held-out six-month test set ($n=279$), XGBoost achieved a \gls{PR-AUC} of 0.456 (95\% \gls{CI}: $[0.364,\,0.568]$), against a rapid-progression prevalence of 0.301. Random Forest obtained 0.459, but the test partition was not used to reopen model selection. At the development-derived threshold of $t=0.21$, XGBoost detected 72 of 84 rapid progressors, giving recall of 85.7\%, precision of 35.1\%, $F_1=0.498$, and $F_2=0.665$. This operating point favoured sensitivity: 133 of the 205 positive predictions were false positives. The Brier score of 0.197 was slightly below the constant-prevalence reference of approximately 0.210, but the reliability plot did not support a claim of full calibration.

The core baseline block retained most of the cross-validated discrimination. Adding \gls{FVC} produced the clearest gain for XGBoost and Random Forest, while vital signs, treatment variables, and laboratory measurements had smaller and less consistent effects. Baseline \gls{ALSFRS-R} also carried the largest global \gls{SHAP} attribution and the largest absolute Logistic Regression coefficient. FVC, respiratory rate, and age recurred across model-level and local explanations, although \gls{SHAP} and \gls{LIME} shared only 50--70\% of their ten leading features in the four case studies. Error analysis added an important qualification: recall fell to 54.5\% among participants with baseline \gls{ALSFRS-R} above 42, the subgroup with the most preserved function.

Categorical-aware LIME preserves the original classifier probabilities but does not explain them exactly. Its local weighted $R^2$ ranges from 0.198 to 0.456; for the false negative, the surrogate even places the probability above the selected threshold. Correct feature representation and partial SHAP agreement therefore do not establish a faithful local explanation.

The supplementary analyses did not change the primary model or threshold. The decision curve's point estimates favoured the model across $t=0.18$--$0.47$, but the paired bootstrap interval for its net-benefit gain over treat all at $t=0.21$ includes zero ($[-0.015,\,0.043]$). No clinical intervention was defined. Training-size sensitivity showed modest average gains without a clear plateau. The patient report demonstrated prediction and explanation generation, not suitability for clinical decisions.

% ============================================================
\section{Main Findings from Study~II}
\label{sec:conclusion_study2}
% ============================================================

The final Study~II cohort contained 1,502 patients, including 174 Short Survivors. Seventy classifier--strategy combinations were tuned on the training partition, and one selected configuration per classifier was subsequently fitted for held-out evaluation. A separate paired analysis examined BalancedBagging while keeping each base classifier's hyperparameters fixed.

BalancedBagging increased mean cross-validated balanced accuracy for all seven classifiers; the difference was statistically significant for six after Bonferroni correction. Among the five families directly comparable with Papaiz et al.~\cite{Papaiz2024}, four showed significant gains, while the small Naïve Bayes increase was not significant after correction. This reproduces the qualitative pattern of the reference study but remains a partial replication: the cohorts, optimisation objectives, and ensemble implementation differed, Random Forest used a different balanced ensemble in the reference work, and the repeated folds were not independent experimental replications.
The additional fixed-configuration comparison shows why this claim should remain specific to the operating point. BalancedBagging improves ROC-AUC for KNN and Decision Tree, leaves Naive Bayes almost unchanged, and lowers it for the other four families. In SVM, balanced accuracy rises from 0.500 to 0.793 despite a ROC-AUC decrease from 0.910 to 0.880.

LightGBM with random oversampling gave the strongest balance at the default threshold on the held-out test set ($n=301$, 35 Short Survivors). Its \gls{ROC-AUC} was 0.923 (95\% \gls{CI}: $[0.868,\,0.966]$), with sensitivity of 0.857 and specificity of 0.831. SVM with SMOTEENN detected one additional Short Survivor, reaching sensitivity of 0.886, but its specificity was lower at 0.767. A paired sensitivity comparison gives exact $p=1.000$ and no clear separation. LightGBM's average precision is 0.704, but its Brier score of 0.1034 is slightly worse than the constant-prevalence reference of 0.1028; its probabilities are not calibrated clinical risks.

The MLP achieved cross-validated and test \gls{ROC-AUC} values of 0.915 and 0.918, yet detected only one of 35 Short Survivors at the default threshold. A sensitivity-constrained threshold selected from out-of-fold training predictions increased test sensitivity to 0.771 while retaining specificity of 0.883.

Diagnosis delay had the largest \gls{SHAP} attribution in both LightGBM and MLP, followed by different combinations of \gls{ALSFRS-R} item slopes. Agreement across the full rankings was moderate ($\rho=0.389$, $p=0.066$). The entry-time audit adds a more specific caution: the median onset-to-trial-reference interval is 11.2 months for Short Survivors and 26.1 for Standard Survival. Late entry partly constrains the onset-based class, so diagnostic delay's prominence cannot be attributed only to disease tempo.

The temporal sensitivity analyses materially narrow the original performance claim. Under an entry-time limit of 18 months, LightGBM+ROS reaches ROC-AUC 0.828 and MLP 0.823, rather than 0.923 and 0.918. A time-bounded 18-month landmark leaves 161 complete cases and only five subsequent deaths in TEST; the corresponding ROC-AUC values are 0.693 and 0.764. This is conditional prediction for months 18--24, not the original retrospective task. The results show sensitivity to timing and population selection, but the altered cohorts and prediction question prevent attributing the entire decrease to left truncation alone.

% ============================================================
\section{Cross-Study Methodological Lessons}
\label{sec:conclusion_lessons}
% ============================================================

The two studies do not support one universal modelling protocol. The primary Study~I analysis used a percentile-based progression target, GroupKFold, and \gls{PR-AUC}; repeated stratified validation then checked sensitivity to the partition. Study~II used \gls{ROC-AUC} to separate ranking performance from the operating point, used repeated stratified folds, and compared explicit resampling strategies. Its temporal analyses changed the observed class balance substantially. The common requirement was not identical methodology, but a clear link between the target, development procedure, and reported claim.

Model ranking was less decisive than a single best score might imply. XGBoost led Study~I cross-validation by 0.003 \gls{PR-AUC} over Random Forest, while Random Forest led the held-out comparison by approximately 0.002 using the unrounded estimates. In Study~II, the three leading cross-validation \gls{ROC-AUC} values were separated by 0.004, yet their default-threshold sensitivities differed substantially. Prespecified selection rules prevented the test sets from becoming a second model-development stage, but the narrow margins also argue against describing one algorithm as intrinsically superior.

Threshold choice changed the practical reading of both studies. Study~I's exploratory threshold of 0.21 recovered 85.7\% of rapid progressors at the cost of 133 false positives; no clinical cost ratio was available to justify that trade-off. In Study~II, the MLP's sensitivity changed from 0.029 to 0.771 after applying a development-derived threshold. Ranking metrics answer whether higher-risk patients tend to receive higher scores; they do not specify how many patients will be detected once a decision boundary is imposed.

Reproducibility depended on correcting imperfections as well as reporting them. Study~I kept final target construction, preprocessing, optimisation, and threshold selection outside the test partition. In Study~II, MinMax scaling was moved inside training partitions and cross-validation folds. For the original complete-case cohort, the transformed matrices and held-out metrics remained unchanged because every ordinal predictor had the same range across those partitions. The temporal analyses use their own fold-fitted transformations; the range equality is not assumed to extend to them. Errors are now recorded rather than silently converted to scores, although this cannot recover the original missing fold-level history.

Model explanations served as a diagnostic layer. The recurrence of functional and respiratory variables in Study~I, and of diagnosis delay and item-level decline in Study~II, made the fitted rules clinically recognisable. Differences between \gls{SHAP}, \gls{LIME}, coefficients, and the two Study~II models also showed that an explanation depends on the model and explainer being examined. Plausibility can reveal unexpected behaviour; it cannot substitute for validation or establish a causal mechanism.

% ============================================================
\section{Limitations}
\label{sec:conclusion_limitations}
% ============================================================

Both studies provide internal rather than external validation. Their training and test partitions come from the same \gls{PRO-ACT} extract, which pools selected clinical-trial populations and heterogeneous recording protocols. Patient separation prevents direct reuse of the same records, but it cannot show that either model transfers to routine care, another registry, or a later recruitment period.
Source trials are not separated either. Predictors and missingness can reflect protocols shared by training and test patients, making patient-level holdout performance potentially optimistic for an unseen trial. A reliable trial-membership variable was not established, so no leave-trial-out claim is made. Study~I's repeated validation assesses partition sensitivity but reuses parameters selected under the original folds; it does not independently evaluate the full selection procedure.

Missingness affects the two studies differently. In Study~I, baseline \gls{FVC} in litres is absent for 43.7\% of the six-month cohort and is handled through fold-fitted median imputation. A fixed-model sensitivity analysis gave development PR-AUC 0.426 with five iterative imputations, compared with 0.434 under the primary strategy; the corresponding post-hoc TEST values were 0.499 and 0.456. In Study~II, only 1,502 of 5,802 eligible patients had all predictors. Expanding the cohort through iterative imputation gave ROC-AUC 0.906 and PR-AUC 0.523 on a different test partition, compared with 0.923 and 0.704 in the complete-case analysis. The lower PR-AUC reinforces the sensitivity of the estimate to cohort construction. Neither analysis proves \gls{MAR} or excludes \gls{MNAR} and selection bias.

Study~I's sparse race indicators also depend on recording conventions. Distinguishing explicit zero, explicit one and unrecorded values, or omitting the fields, changes mean development PR-AUC by less than 0.003 for fixed XGBoost and LR configurations (Appendix~\ref{app:review_s1_race}). This limited ranking sensitivity does not validate the original median coding or establish subgroup fairness.

The outcome definitions also restrict interpretation. Study~I's 30th-percentile cutoff is an operational definition of rapid progression, not a clinical boundary, and eligibility at six months may favour patients with sufficient follow-up. Study~II converts survival into a fixed 24-month class and excludes patients whose status at that horizon is unresolved. Its closest-record rule also yielded post-diagnosis ALSFRS, FVC, and weight measurements for more than 99\% of the final cohort. Neither representation captures the complete longitudinal disease course, and Study~II cannot support diagnosis-time use.
Study~II additionally has left truncation: patients must survive until trial entry, and 756 Standard Survival participants are already at or beyond 24 months at the source-trial reference day. Entry restriction tests sensitivity to that selection, while landmarking defines a common prediction time among observed survivors. Neither estimates unconditional prognosis from onset or removes trial-recruitment and complete-case selection. The 12-month landmark has too few negative cases for the planned validation, and the 18-month landmark has only five positive test outcomes. Their limitations cannot be repaired by reporting the original cohort's larger sample size.

The held-out estimates remain imprecise. Study~I has 84 rapid progressors in its test partition and a \gls{PR-AUC} interval of $[0.364,\,0.568]$. Study~II has only 35 Short Survivors; one additional correct prediction changes sensitivity by 2.86 percentage points, and the LightGBM sensitivity interval spans $[0.724,\,0.967]$. Bootstrap intervals condition on the fitted model and selected threshold, while the repeated cross-validation tests reuse overlapping training observations. They do not represent the variation expected if cohort construction and model development were repeated from the start.

Finally, no prospective clinical use was evaluated. The Study~I report is a software demonstrator, and the decision curve does not correspond to a specified intervention. Study~II predicts status relative to an onset-based boundary rather than remaining survival from diagnosis. Neither output should be used for treatment selection or prognostic counselling on the evidence available here.

% ============================================================
\section{Future Work}
\label{sec:conclusion_future}
% ============================================================

External validation is the first requirement for both pipelines. Study~I should be evaluated in population-based ALS registries and multicentre routine-care cohorts outside PRO-ACT, with the target, threshold, and evaluation plan fixed in advance. Its most plausible initial users would be trial investigators assessing cohort stratification and multidisciplinary ALS teams considering whether a patient may require closer follow-up, rather than patients receiving an automated prognosis. Study~II needs a larger cohort with dated predictors, a fixed prediction time, and enough subsequent events to develop and evaluate conditional risk. The landmark analysis establishes a clearer temporal task, but its 33 test patients cannot support clinical translation. Calibration and subgroup behaviour should be reassessed in independent data rather than assumed to transfer with discrimination.

A prospective evaluation should begin with both models locked and operating silently across several clinical centres. Predictions would be recorded without being shown to clinicians, allowing discrimination, calibration, subgroup performance, missing-data behaviour, and workflow feasibility to be assessed under routine conditions. Only after acceptable external and prospective performance should a model-assisted study compare usual care with access to the prediction. Such a study would need to define the decision supported by each risk estimate, the relative consequences of false negatives and false positives, and whether model use changes trial recruitment, follow-up planning, or patient outcomes. This sequence separates technical validation from the later question of clinical benefit.

The missing-data sensitivity analyses should be repeated in an external cohort under a prespecified protocol. Study~I would benefit from imputers designed explicitly for mixed continuous and categorical data, models with native missing-value handling, and checks on the stability of feature attributions. Study~II requires investigation of why predictor availability differs and whether the expanded imputed cohort is closer to the intended population. Pattern-mixture or other \gls{MNAR}-oriented analyses would test assumptions that chained equations cannot verify. A dynamic progression model could also update risk as new \gls{ALSFRS-R} and respiratory measurements arrive instead of compressing follow-up into one baseline vector and one final class.

For Study~II, a time-to-event analysis with delayed entry would retain censored observations while accounting for survival up to trial inclusion. The fitted entry-restricted and landmark analyses provide starting points, but need extension beyond the current selected complete cases. Both studies also need validation that separates source trials rather than only patients. Future protocols should enforce predictor availability, assess calibration after resampling, and select thresholds against a prespecified clinical objective. TRIPOD+AI~\cite{collins2024tripodai} should guide prospective reporting; the retrospective coverage map in Appendix~\ref{app:tripod_checklist} records what this dissertation reports and what remains unassessed.

The patient report should only progress after predictive validation. Its next evaluation should involve clinicians and patients in deciding what probability, uncertainty, explanation, and warning information can be understood without overstating the model. The same principle applies to SHAP comparisons: stability across resamples, cohorts, and model families is more informative than adding further isolated explanation plots.

The replication framework used in Study~II can be applied to other published \gls{ALS} prognosis pipelines. Reconstructing their targets and cohorts, separating reproduced findings from deliberate extensions, and reporting threshold-dependent errors would help determine which conclusions persist beyond one implementation.

% ============================================================
\section{Final Remarks}
\label{sec:conclusion_final}
% ============================================================

Study~I produced a separated evaluation of rapid functional progression, with held-out \gls{PR-AUC} of 0.456 and recall of 85.7\% at its selected threshold; repeated validation did not preserve its original model ordering. Study~II reproduced the qualitative BalancedBagging pattern and obtained retrospective reference ROC-AUC 0.923 with LightGBM+ROS, but discrimination was lower under temporal restrictions. Neither study identifies a universally best classifier or establishes readiness for clinical deployment.

Explainable and reproducible \gls{ALS} prognosis modelling is possible with public clinical-trial data, but credibility rests on the analysis surrounding the model: a defensible target, strict separation of development and evaluation, imbalance-aware metrics, explicit threshold selection, uncertainty reporting, and explanations treated as descriptions of fitted behaviour. The two studies establish that workflow and its limitations. External and prospective evaluation remain necessary before either pipeline can support individual care.
